% !TeX root = ../../main.tex
\subsection{Descendant and volume coefficients}
\label{subsec:a:01}

We use the normalized intersection recursions of
\cite[Section~2]{MZ15}, in the form recalled in
\cite[Theorem~4.3]{HMT25}. The factorial-normalized brackets are those
of~\eqref{eq:s03:descendants}. To identify their first two coefficients
we use \cite[Lemma~7.2 and Appendix~C.1]{W26}; its common amplitude
factor at the WP specialization cancels in the quotient by $V_{G,n}$.
Throughout this subsection put $k=2G-2+n$. For a nonnegative integer
tuple $\mathbf d=(d_1,\ldots,d_n)$ put
\[
 D=\sum_{i=1}^n d_i,\qquad
 m_a=\#\{i:d_i=a\}\quad(a=0,1,2).
\]
Thus the number of positive entries is $n-m_0$. Define the signed
coefficients $e_{n,r}(\mathbf d)$, for integers $r\geq1$, in the
unique full fixed-tuple expansion supplied by HMT. Its first terms are
\begin{equation}\label{eq:a:profile-definition}
 \frac{[\tau_{\mathbf d}]_{G,n}}{V_{G,n}}
 =1+\frac{e_{n,1}(\mathbf d)}k
     +\frac{e_{n,2}(\mathbf d)}{k^2}+O_{n,\mathbf d}(k^{-3}).
\end{equation}
For the empty tuple both coefficients are zero. The convention also
applies to the zero tuple of every positive length.

We first give a closed potential for W26's discrete second-coefficient
recursion. The low-degree counts in this potential will subsequently be
summed, rather than discarded as exceptional modes.

\begin{prp}[The first two profiles]\label{prp:a:profiles}
For every fixed $n\geq1$ and $\mathbf d\in\mathbb Z_{\geq0}^n$,
form the following local polynomials in its statistics:
\begin{align*}
 U={}&24D^2+24Dn-50D+m_0^3-19m_0-n^3+19n+6m_1,\\
 V={}&80D^3+(120n-372)D^2+(60n^2-372n+754)D\\
 &+8m_0^3+(36D+18m_1+18n-141)m_0^2\\
 &+(-180D-90m_1-90n+547)m_0\\
 &-26n^3+231n^2-547n+102m_1+30m_2.
\end{align*}
Then
\begin{align}
 e_{n,1}(\mathbf d)
 &=-\frac{2D^2+(2n-5)D+(m_0^2-5m_0-n^2+5n)/2}{\pi^2},
 \label{eq:a:first-profile}\\
 e_{n,2}(\mathbf d)
 &=\frac{e_{n,1}(\mathbf d)^2}{2}+\frac{n-m_0}{16}
                  -\frac{U}{12\pi^2}-\frac{V}{12\pi^4}.
 \label{eq:a:second-profile}
\end{align}
In an expansion in $G^{-1}$ the second coefficient is
\begin{equation}\label{eq:a:profile-genus-conversion}
 \frac{[\tau_{\mathbf d}]_{G,n}}{V_{G,n}}
 =1+\frac{e_{n,1}(\mathbf d)}{2G}
 +\frac{e_{n,2}(\mathbf d)-(n-2)e_{n,1}(\mathbf d)}{4G^2}
 +O_{n,\mathbf d}(G^{-3}).
\end{equation}
\end{prp}
\begin{prf}
The first formula is W26's deficit with its sign reversed. Here is a
finite verification of the second, including all pivot degrees. Define
recursion coefficients $\mu_l$ and their tails by
\begin{equation}\label{eq:a:recursion-kernel}
 \sum_{l\geq0}\mu_l z^l
 =(1-z)\frac{\pi\sqrt z}{\sin(\pi\sqrt z)},\qquad
 T_j=\sum_{l\geq\max(0,j)}\mu_l\quad(j\in\mathbb Z).
\end{equation}
The quotient is interpreted by its power series at zero. In particular
$\mu_0=1$, $\mu_1=\pi^2/6-1$, and
$\mu_2=7\pi^4/360-\pi^2/6$.

Fix a pivot $i$, write $t=d_i$, and let $z_0=m_0-\mathbf1_{t=0}$ be
the number of zero entries away from it. The symbol $\mathbf1$ denotes
an indicator. Let $\mathbf d^{ij,l}$ merge entries $i,j$ into
$t+d_j+l-1$, and let $\mathbf d^{i;u,v}$ replace entry $i$ by $u,v$.
Only nonnegative resulting entries are included, and $\mathbf e_i$
denotes the $i$th unit vector. W26's (C.1)--(C.3)
give the increment $I_i=e_{n,2}(\mathbf d+\mathbf e_i)-e_{n,2}(\mathbf d)$
as
\begin{align}
 I_i={}&-\frac1{\pi^2}\sum_{j\ne i}(2d_j+1)
  \sum_{l\geq\max(0,1-t-d_j)}\mu_l
  \left[\frac32+\frac{2n-6}{\pi^2}
                 +e_{n-1,1}(\mathbf d^{ij,l})\right]\notag\\
 &-\frac2{\pi^2}\sum_{l\geq\max(0,2-t)}\mu_l
  \sum_{\substack{u+v=t+l-2\\u,v\geq0}}
  \left[\frac32+\frac{4-2n}{\pi^2}
                 +e_{n+1,1}(\mathbf d^{i;u,v})\right]\notag\\
 &-\frac1{\pi^4}
 \left[\left(\binom{z_0}{2}+\frac{\pi^2}{12}\right)T_{2-t}
                                      +\frac12 T_{3-t}\right].
 \label{eq:a:second-increment}
\end{align}
The merging sum is absent for $n=1$.

We reduce the increment to four cases. For $t\geq3$ its value is
\begin{align*}
 I_*={}&-\frac{4D+2n-13/6}{\pi^2}+\frac1{\pi^4}
 \bigl[8D^3+(12n-38)D^2\\
 &+(2m_0^2-10m_0+2n^2-18n+45)D
 +(n-9/2)m_0^2\\
 &+(-5n+45/2)m_0-n^3+7n^2/2+15n/2-34\bigr].
\end{align*}
The two positive low modes satisfy
\[
 I_i-I_*=\frac{\pi^2+3m_0^2-15m_0+12}{2\pi^4}\ (t=1),\qquad
 I_i-I_*=\frac5{2\pi^4}\ (t=2).
\]
For the zero mode the reduction is
\begin{align*}
 16\pi^4(I_i-I_*)={}&(-32m_0+96)D^2
 +(-32m_0n+112m_0+96n-336)D\\
 &+\pi^4+(4m_0^2-4m_0-32)\pi^2
 +(48m_0-144)m_1\\
 &-8m_0^3+80m_0^2+(8n^2-24n-264)m_0\\
 &-24n^2+72n+288\qquad(t=0).
\end{align*}
To obtain these four expressions one needs only the moments
$\sum\mu_l l^j$ for $0\leq j\leq3$ in the table below. The ordered
pair sum has two endpoint pairs and $N-1$ interior pairs when its total
degree $N\geq1$; its $N=0$ value is evaluated separately.

On the proposed potential, incrementing the pivot changes the statistics
by the following amounts:
\[
\begin{array}{c|rrrr}
 t&\Delta D&\Delta m_0&\Delta m_1&\Delta m_2\\\hline
 0&1&-1&1&0\\
 1&1&0&-1&1\\
 2&1&0&0&-1\\
 t\geq3&1&0&0&0
\end{array}
\]
Substitution in~\eqref{eq:a:second-profile} gives $I_i$ in each of
the four cases. This verifies every coordinate increment, rather
than fitting values at a collection of tuples. The potential vanishes
at $D=0,m_0=n,m_1=m_2=0$; the increment recursion therefore determines
it uniquely on all nonnegative tuples. Finally expand
$k^{-1}=(2G)^{-1}-(n-2)(2G)^{-2}+O(G^{-3})$.
\end{prf}

\begin{rmk}[Identification]
The source quantities in W26 are evaluated at its parameter $Y=0$ and
complexity $2G-2+n$. Its deficit is $-e_{n,1}$, and its second coefficient
is $e_{n,2}$. Lemma~C.1 itself fixes the number of markings. Uniqueness
of each fixed-tuple expansion identifies these formulas with the HMT
coefficients in~\eqref{eq:s03:hmt-descendant}. The polynomial remainders
and summable degree bounds proved in Section~\ref{sec:s03}, rather than
an unstated uniform version of W26, justify their later use when the
loop count grows. The $m_1,m_2$ terms remain part of those coefficients.
\end{rmk}

The absolute second cusp-volume coefficient requires one more order of
an exact recursion. We give the extraction rule explicitly before its
evaluation. All auxiliary ratios and coefficients in the rest of this
subsection are local to this calculation.

Define the signed dilaton weights
\begin{equation}\label{eq:a:dilaton-weights}
 \nu_d=\frac{2(-1)^{d-1}d\pi^{2d}}{(2d+1)!}\quad(d\geq1),\qquad
 \sum_{d\geq1}\nu_d z^d
   =f(2\pi^2z)-\cos(\pi\sqrt z).
\end{equation}
The kernel and dilaton moments needed below are finite evaluations of
their generating functions at $z=1$:
\begin{equation}\label{eq:a:finite-moments}
\begin{array}{c|c|c}
 j&\sum_{l\geq0}\mu_l l^j&\sum_{d\geq1}\nu_d d^j\\\hline
0&2&1\\
1&3/2&-1/2\\
2&\pi^2/6+5/4&(1-\pi^2)/4\\
3&5\pi^2/8+9/8&-\pi^2/4-1/8\\
4&7\pi^4/120+5\pi^2/3+17/16&(\pi^4-5\pi^2+1)/16\\
5&49\pi^4/96+125\pi^2/32+33/32&(9\pi^4-10\pi^2-1)/32\\
6&\text{not needed}&(-\pi^6+56\pi^4-21\pi^2+1)/64
\end{array}
\end{equation}
For example, with the local variable $r=\sqrt z-1$, the two generating
functions become
\[
 (2+3r+r^2)\frac{\pi r}{\sin(\pi r)},\qquad
 \cos(\pi r)-\frac{\sin(\pi r)}{\pi(1+r)}.
\]
Applying $(z\partial_z)^j$ proves every entry of the table.

Write the two exact volume ratios as
\begin{equation}\label{eq:a:local-ratios}
 \mathcal C_n(k)=4\pi^2 k\frac{V_{G,n}}{V_{G,n+1}},\qquad
 \mathcal V_n(k)=\frac{V_{G-1,n+2}}{V_{G,n}},
 \qquad k=2G-2+n,
\end{equation}
and denote their coefficients of $k^{-r}$ by $u_{r,n}$ and $v_{r,n}$.
In the remaining rules write $e_{n,r}(a)$ for
$e_{n,r}(a,0^{n-1})$ and $e_{n,r}(a,b)$ for
$e_{n,r}(a,b,0^{n-2})$, where the latter requires $n\geq2$.
An overbar means the signed single-degree sum
$\overline{e_{n,r}}=\sum_{d\geq1}\nu_d e_{n,r}(d)$.
The exact dilaton identity is
\begin{equation}\label{eq:a:exact-cusp}
 \mathcal C_n(k)
 =\sum_{d\geq1}\nu_d
       \frac{[\tau_d\tau_0^n]_{G,n+1}}{V_{G,n+1}}.
\end{equation}
Dimensionally inadmissible brackets are zero. Its child complexity is
$k+1$, so, with an order-three profile coefficient supplied by HMT,
\begin{align}
 u_{1,n}&=\overline{e_{n+1,1}}
                 =\frac12+\frac{2n-4}{\pi^2},\notag\\
 u_{2,n}&=-\overline{e_{n+1,1}}+\overline{e_{n+1,2}},\notag\\
 u_{3,n}&=\overline{e_{n+1,1}}-2\overline{e_{n+1,2}}
                                      +\overline{e_{n+1,3}}.
 \label{eq:a:cusp-extraction}
\end{align}

Here are the exact recursions that determine the remaining coefficients.
For $n\geq2$ the genus identity is
\begin{equation}\label{eq:a:exact-genus}
 \frac{[\tau_1\tau_0^{n-1}]_{G,n}}{V_{G,n}}
 =\mathcal V_n(k)+\frac6{V_{G,n}}
 \sum_{\substack{G_1+G_2=G\\I\sqcup J=\{1,\ldots,n-2\}}}
        V_{G_1,|I|+2}V_{G_2,|J|+2}.
\end{equation}
Only stable factors are included. To state the adjacent-degree identity,
take the first entry $t$ of a tuple as pivot. Subtract the MZ recursion
at $t+1$ from that at $t$, and shift its kernel index. The result is
\begin{align}
 &[\tau_{\mathbf d}]_{G,n}
       -[\tau_{\mathbf d+\mathbf e_1}]_{G,n}\notag\\
 &\quad=4\sum_{j=2}^n(2d_j+1)\sum_{l\geq0}\mu_l
                            [\tau_{\mathbf d^{1j,l}}]_{G,n-1}\notag\\
 &\qquad+8\sum_{l\geq0}\mu_l
       \sum_{\substack{u+v=t+l-2\\u,v\geq0}}
                   [\tau_u\tau_v\prod_{i=2}^n\tau_{d_i}]_{G-1,n+1}\notag\\
 &\qquad+8\sum_{l\geq0}\mu_l
       \sum_{\substack{u+v=t+l-2\\u,v\geq0}}
       \sum_{\substack{G_1+G_2=G\\I\sqcup J=\{2,\ldots,n\}}}
       [\tau_u\prod_{i\in I}\tau_{d_i}]_{G_1,|I|+1}
       [\tau_v\prod_{i\in J}\tau_{d_i}]_{G_2,|J|+1}.
 \label{eq:a:exact-adjacent}
\end{align}
Negative degrees and unstable children contribute zero. This identity
is used for large $G$ with fixed $n$; the exceptional initial parent
types are not needed. The factors $4,8,8$ use
$\mu_l=2(a_l-a_{l-1})$ in the MZ/HMT kernel notation. The last sum is
ordered, so choosing the smaller component combines two orientations.

We next justify extracting a coefficient from these infinite recursions.
This also controls components whose genus grows with the parent genus.

\begin{lmm}[Remainders in the finite extraction]\label{lmm:a:extraction-tail}
For fixed marking numbers and fixed pivot degree, extracting
\eqref{eq:a:exact-adjacent} through order $k^{-3}$ has remainder
$O(k^{-4})$. In its separating sum it suffices to retain components
with loss $2j+s\leq3$, where $j$ is their genus and $s$ the number
of nonpivot labels on them. In~\eqref{eq:a:exact-genus} at $n=2$,
only the $(1,2)$ component contributes through that order.
Polynomially weighted kernel sums and the signed dilaton average
preserve these remainder orders.
\end{lmm}
\begin{prf}
The singularity at $z=1$ in~\eqref{eq:a:recursion-kernel} is removable;
the next one is at $z=4$. Thus $\sum|\mu_l|(1+l)^a<\infty$ for every
fixed nonnegative integer $a$. HMT's order-two descendant and ratio
remainders are polynomial in the child degrees. Their volume prefactors
in~\eqref{eq:a:exact-adjacent} are $O(k^{-1})$, so their summed errors
are $O(k^{-4})$. Terms beyond a dimension cutoff have exponentially
small kernel weight. The HMT order-three expansion likewise supplies
the actual third profile and its polynomial majorant. Finally
$\sum|\nu_d|(1+d)^a<\infty$ at every fixed $a$.

We bound the full omitted genus sum by factorials. If the smaller child
has complexity $b$ in the adjacent identity,
the other has complexity $k-1-b$. The all-type volume bound and the
fixed-marking lower asymptotic give
\[
 \frac{V_{j,s+1}V_{G-j,n-s}}{V_{G,n}}
 \leq C_n\frac{\Gamma(b)\Gamma(k-1-b)}{\Gamma(k)}.
\]
The square-root genus factor is bounded because the other child has
at least half the total complexity. There are finitely many label
subsets, and normalized descendant brackets are at most one.
The omitted losses have $b\geq3$. Put $N=k-1$; the consecutive ratio
of $\Gamma(b)\Gamma(N-b)$ is $b/(N-b-1)$. Up to $N/4$ it is at
most $1/2$ for large $N$, and from $N/4$ to $N/2$ the terms decrease.
Consequently
\[
 \sum_{3\leq b\leq(k-1)/2}
 \frac{\Gamma(b)\Gamma(k-1-b)}{\Gamma(k)}=O(k^{-4}).
\]
More explicitly, the interval $N/4\leq b\leq N/2$ is bounded by
$N$ times its term at $\lfloor N/4\rfloor$, which is exponentially
small relative to the term at $b=3$. The ordered node-degree count
only adds a polynomial kernel moment.
For loss two use a first profile correction, and for loss three use
the leading profile; their errors also have order $k^{-4}$.

In the genus identity the complexities sum to $k$. After removing
the $(1,2)$ component the smaller complexity is at least four.
The same ratio argument gives
$\sum_{4\leq b\leq k/2}\Gamma(b)\Gamma(k-b)/\Gamma(k)=O(k^{-4})$.
This proves the asserted error for all separating genera together.
For the dilaton average use HMT's order-three remainder, polynomial
in the single degree, and the factorial tail of $\nu_d$; thus the
averaged error has a uniform summable majorant, rather than a
degree-dependent fixed-tuple constant.
\end{prf}

We describe the finite calculation of the third profile.
For $n\geq2$ the first two genus-ratio coefficients are
\begin{align}
 v_{1,n}&=-\frac{4n-6}{\pi^2},\notag\\
 v_{2,n}&=e_{n,2}(1)
 -\frac{3(n-2)}{\pi^2}(1+u_{1,n-1})
 -\frac{3}{2\pi^2}\binom{n-2}{2}-\frac3{16}.
 \label{eq:a:genus-first-two}
\end{align}
These follow by retaining $(0,3),(0,4),(1,2)$ in
\eqref{eq:a:exact-genus}. For $n=2$ the precise retained term is
\begin{equation}\label{eq:a:genus-small-split}
 \frac{12V_{1,2}V_{G-1,2}}{V_{G,2}}
   =\frac3{16k^2}+\frac3{4k^3}+O(k^{-4}),\qquad
 v_{3,2}=e_{2,3}(1)-\frac34.
\end{equation}
The value $V_{1,2}=\pi^4/4$ is derived in
Subsection~\ref{subsec:a:03} without using any asymptotic coefficient.

For the one-positive tuple put
\[
 M_1=1+u_{1,n-1},\quad M_2=1+2u_{1,n-1}+u_{2,n-1},
\]
\[
 N_1=1+u_{1,n+1}+v_{1,n},\quad
 N_2=1+2u_{1,n+1}+u_{2,n+1}+v_{2,n}
                              +v_{1,n}(1+u_{1,n+1}).
\]
They give the prefactor expansions
\begin{align*}
 4\frac{V_{G,n-1}}{V_{G,n}}
 &=\frac1{\pi^2k}\left(1+\frac{M_1}{k}+\frac{M_2}{k^2}
                                      +O(k^{-3})\right),\\
 8\frac{V_{G-1,n+1}}{V_{G,n}}
 &=\frac2{\pi^2k}\left(1+\frac{N_1}{k}+\frac{N_2}{k^2}
                                      +O(k^{-3})\right).
\end{align*}
Both child complexities are $k-1$.

Let $\mathcal D_n(t)$ denote the coefficient of $k^{-3}$ in the normalized
adjacent difference at $(t,0^{n-1})$. Its merging and nonseparating
parts are, respectively,
\begin{align}
 &\frac{n-1}{\pi^2}\sum_{l\geq0}\mu_l
 \left[M_2+(M_1+1)e_{n-1,1}(t+l-1)
                            +e_{n-1,2}(t+l-1)\right],\notag\\
 &\frac2{\pi^2}\sum_{l\geq0}\mu_l
 \sum_{\substack{u+v=t+l-2\\u,v\geq0}}
 \left[N_2+(N_1+1)e_{n+1,1}(u,v)+e_{n+1,2}(u,v)\right].
 \label{eq:a:third-merge-split}
\end{align}
An entire bracket with negative degree is zero. The loss-two separating
parts, with the same convention, are
\begin{align}
 &\frac{\binom{n-1}{2}}{\pi^4}\sum_l\mu_l
 \left[4+\frac{4n-14}{\pi^2}+e_{n-2,1}(t+l-2)\right],\notag\\
 &\frac1{12\pi^2}\sum_l\mu_l[4+e_{n,1}(t+l-2)],\qquad
 \frac1{2\pi^4}\sum_l\mu_l[4+e_{n,1}(t+l-3)].
 \label{eq:a:third-loss-two}
\end{align}
The sphere expression is absent when $n<3$. The loss-three parts are
\begin{align}
 &\binom{n-1}{3}\left(\frac{T_{2-t}}{2\pi^4}
                                      +\frac{3T_{3-t}}{\pi^6}\right),\notag\\
 &(n-1)\left(\frac{T_{2-t}}{16\pi^2}
             +\frac{T_{3-t}}{2\pi^4}+\frac{5T_{4-t}}{2\pi^6}\right).
 \label{eq:a:third-loss-three}
\end{align}
Adding these expressions gives $\mathcal D_n(t)$, and hence
\begin{equation}\label{eq:a:third-primitive}
 e_{n,3}(d,0^{n-1})=-\sum_{t=0}^{d-1}\mathcal D_n(t).
\end{equation}
The sphere term in~\eqref{eq:a:third-loss-two} has relative volume
correction $4+(4n-14)/\pi^2$. For the torus it is $4$. Multiply the
small-type brackets in~\eqref{eq:a:small-amplitudes} by $16$ for both
orientations to obtain the displayed coefficients. At loss three the
same rule gives~\eqref{eq:a:third-loss-three}. Every small type and
every small node degree is retained.

These are finite rules: in each ordered pair sum the two endpoint pairs
have one positive entry, and the other $a-1$ pairs have two when their
total degree is $a\geq3$. There are two entries equal to one and two
equal to two across these interior pairs. Thus the pair sum of
$e_{n,2}$ is its one-positive value twice, plus $a-1$ times its
two-positive value with $m_1=m_2=0$, plus twice each of the $m_1,m_2$
coefficients. At $a=0,1,2$ use the actual tuples. The analogous rule
for $e_{n,1}$ has no $m_1,m_2$ terms. It follows that $\mathcal D_n(t)$ is a
polynomial of degree at most five for $t\geq5$, with only the five
corrections $t=0,\ldots,4$. Sum its polynomial primitive by the usual
power sums and use~\eqref{eq:a:finite-moments}. A correction at $t=j$
is averaged with $\sum_{d>j}\nu_d=1-\sum_{d=1}^j\nu_d$.
This specifies the calculation without truncating an infinite sum.

\begin{prp}[Second volume coefficient]\label{prp:a:c-two}
The fixed-marking coefficients in~\eqref{eq:s03:fixed-n-volume} satisfy
\begin{align}
 c_2(0)&=\frac{145}{288}-\frac{395}{72\pi^2}
                                  +\frac{1513}{72\pi^4},
 \label{eq:a:c-two-zero}\\
 c_2(n)-c_2(0)&=\frac{3n^2}{32}-\frac{29n}{64}
 +\frac{17n^3/48-67n^2/24+7n}{\pi^2}\notag\\
 &\hspace{15mm}
 +\frac{n^4/8-2n^3+271n^2/24-79n/3}{\pi^4}
 \qquad(n\geq0).
 \label{eq:a:c-two-n}
\end{align}
\end{prp}
\begin{prf}
The preceding rules give the following three expansions in each ratio's
own parent complexity. They are displayed through the order needed to
identify an absolute second coefficient:
\begin{align}
 \mathcal C_2(k)={}&1+\frac1{2k}
 +\frac{5\pi^4+36\pi^2-48}{48\pi^4k^2}\notag\\
 &-\frac{11\pi^6-138\pi^4+792\pi^2-2016}{96\pi^6k^3}
 +O(k^{-4}),\notag\\
 \mathcal C_3(k)={}&1+\frac{\pi^2+4}{2\pi^2k}
 +\frac{5\pi^4+108\pi^2+48}{48\pi^4k^2}\notag\\
 &-\frac{41\pi^6-324\pi^4-336\pi^2-4608}{288\pi^6k^3}
 +O(k^{-4}),\notag\\
 \mathcal V_2(k)={}&1-\frac2{\pi^2k}
 -\frac{3\pi^4+80\pi^2-96}{24\pi^4k^2}\notag\\
 &-\frac{86\pi^6+663\pi^4-2052\pi^2+6192}{144\pi^6k^3}
 +O(k^{-4}).
 \label{eq:a:three-ratios}
\end{align}
Each entry is obtained by the two finite moment columns in
\eqref{eq:a:finite-moments}, the low-mode corrections just specified,
and~\eqref{eq:a:cusp-extraction} or
\eqref{eq:a:genus-small-split}. Lemma~\ref{lmm:a:extraction-tail}
shows that these are coefficients of actual volume ratios.

Put $k=2g$. The exact telescoping identity is
\begin{equation}\label{eq:a:adjacent-genus-match}
 \frac{V_{g+1,2}}{(4\pi^2)^2 k(k+1)V_{g,2}}
 =\frac1{\mathcal C_2(k)\mathcal C_3(k+1)\mathcal V_2(k+2)}.
\end{equation}
The three nonconstant coefficients on its right are
\[
 -1,\qquad \frac{7\pi^2-4}{6\pi^2},\qquad
                    -\frac{3\pi^4+2\pi^2-24}{3\pi^4}.
\]
Apply the actual fixed-$n$ MZ expansion through order three, with
remainder $O(g^{-4})$. Its unknown $c_3(2)$ cancels in this adjacent
ratio. The coefficient of $k^{-3}$ on the left is therefore
\[
 -\frac52+12c_1(2)+8c_1(2)^2-16c_2(2).
\]
The first two coefficients check $-1$ and $3/2-4c_1(2)$ independently.
The third determines
$c_2(2)=-(2\pi^4-13\pi^2+35)/(72\pi^4)$.

Finally the cusp ratio through second order gives
\[
 c_2(n+1)-c_2(n)=-\frac{u_{2,n}}4
 -\left(\frac{2-n}{2}+c_1(n+1)\right)
                         (c_1(n)-c_1(n+1)).
\]
Substitute~\eqref{eq:a:second-profile} into
\eqref{eq:a:cusp-extraction} and use the dilaton moment table.
For $n=0,1$ the increments are respectively
\[
 -\frac{69\pi^4-876\pi^2+3248}{192\pi^4},\qquad
 -\frac{33\pi^4-212\pi^2+880}{192\pi^4}.
\]
Subtracting them from $c_2(2)$ proves~\eqref{eq:a:c-two-zero}.
Summing the same polynomial increment from $0$ to $n-1$ proves
\eqref{eq:a:c-two-n}. No absolute second coefficient was used in
constructing the three ratio rows.
\end{prf}

\begin{rmk}[Prediction]
See \cite[(6.10)]{EGGLS24} for the earlier prediction, which our answer
matches.
The resurgence assumption appears in their Section~6.1 before (6.6).
We use that result for credit and a comparison after our calculation.
The unconditional inputs here
are the MZ/HMT recursions, W26's profiles, and the uniform bounds from
Section~\ref{sec:s03}.
\end{rmk}
